\begin{thema}{Γ}
Ένας κώνος είναι εγγεγραμμένος σε μια σφαίρα σταθερής ακτίνας $R=3$. Έστω $x$ η απόσταση του κέντρου της σφαίρας από τη βάση του κώνου, με $x \in (0, 3)$. 
\begin{center}
\begin{tikzpicture}
    % ==========================================
    % Ορισμός Παραμέτρων 
    % ==========================================
    \def\R{3}       % Ακτίνα σφαίρας R=3
    \def\d{1}       % Απόσταση x 
    \def\r{2.828}   % Ακτίνα βάσης
    \def\yell{0.25} % Προοπτική ελλείψεων 

    % ==========================================
    % 1. Σχεδιασμός Σφαίρας (με Radial Shading)
    % ==========================================
    \shadedraw[shading=radial, inner color=white, outer color=gray!30, draw=gray!80, thick] (0,0) circle (\R);
    
    % Προαιρετικά: Προσθέτουμε ένα "γυάλισμα" 
    \fill[white, opacity=0.3] (-0.3*\R, 0.4*\R) ellipse ({0.2*\R} and {0.15*\R});

    % Ισημερινός σφαίρας 
    \draw[dashed, color=gray!60] (\R,0) arc (0:180:{\R} and {\R*\yell});
    \draw[color=gray!60] (-\R,0) arc (180:360:{\R} and {\R*\yell});

    % ==========================================
    % 2. Σχεδιασμός Βάσης Κώνου 
    % ==========================================
    \coordinate (K) at (0, -\d);
    
    \shadedraw[dashed, thick, color=\xrwma!60!black, 
               top color=\xrwma!10, bottom color=\xrwma!20, fill opacity=0.6] 
        (\r, -\d) arc (0:180:{\r} and {\r*\yell})
        arc (180:360:{\r} and {\r*\yell});

    % ==========================================
    % 3. Σχεδιασμός Πλευρών Κώνου (με Linear Shading)
    % ==========================================
    \coordinate (V) at (0, \R);               
    \coordinate (Left) at (-\r, -\d);         
    \coordinate (Right) at (\r, -\d);         
    
    % Shading πλευρών
    \shadedraw[thick, color=\xrwma!60!black,
               left color=\xrwma!60, right color=\xrwma!60, middle color=white!80!\xrwma, fill opacity=0.8] 
               (V) -- (Left) arc (180:360:{\r} and {-\r*\yell}) -- cycle;

    % ==========================================
    % 4. Γεωμετρικά Στοιχεία & Ετικέτες
    % ==========================================
    \coordinate (O) at (0, 0);                
    \coordinate (A) at (\r, -\d);             
    
    % Κατακόρυφος άξονας
    \draw[dashed, thick] (V) -- (O);
    \draw[thick, secondcolor] (O) -- (K) node[midway, left] {$x$};
    
    % Ακτίνα βάσης κώνου (ρ)
    \draw[dashed, thick, \xrwma!60!black] (K) -- (A) node[midway, below] {$\rho$};
    
    % Ακτίνα σφαίρας (R)
    \draw[thick, \xrwma!50!black] (O) -- (A) node[midway, above] {$R=3$};

    % Σημεία
    \filldraw (O) circle (1pt) node[left] {$O$};
    \filldraw (K) circle (1pt) node[below left] {$K$};
    \filldraw (V) circle (1pt) node[above] {$V$};
    \filldraw (A) circle (1pt) node[right] {$A$};

\end{tikzpicture}
\end{center}
\begin{erwthma}
  \item Να αποδείξετε ότι η ακτίνα $\rho$ της βάσης του κώνου δίνεται από τη σχέση $\rho^2 = 9 - x^2$.
  \item Αποδείξτε ότι ο όγκος του κώνου συναρτήσει του $x$ δίνεται από τον τύπο $V(x) = \dfrac{\pi}{3}(9-x^2)(3+x)$.
  \item Να βρείτε την τιμή του $x$ για την οποία ο όγκος του κώνου μεγιστοποιείται, καθώς και τις διαστάσεις (ύψος και ακτίνα) του μέγιστου κώνου.
  \item Αν το $x$ αυξάνεται με ρυθμό $0.5 \si{cm/s}$, να βρείτε το ρυθμό μεταβολής του όγκου $V(x)$ τη χρονική στιγμή που το ύψος του κώνου είναι $4\si{cm}$.
\end{erwthma}
\end{thema}
